% ============================================
% Postdoctoral Fellowship Cover Letter – Georgetown University
% Center for Business and Public Policy
% ============================================

\documentclass[11pt,letterpaper]{letter}

\usepackage{geometry}
\usepackage{microtype}
\usepackage[hidelinks]{hyperref}

\geometry{left=25mm,right=25mm,top=25mm,bottom=25mm}

\signature{Decory Edwards}
\address{
Department of Economics \\
Johns Hopkins University \\
Baltimore, MD 21218 \\
\texttt{dedwar65@jh.edu}
}

\begin{document}

\begin{letter}{Search Committee \\
Georgetown Center for Business and Public Policy \\
McDonough School of Business \\
Georgetown University \\
Washington, DC \\ USA}

\opening{Dear Members of the Search Committee,}

I am writing to apply for the Postdoctoral Fellowship in Applied Microeconomics and Industrial Organization at the Georgetown Center for Business and Public Policy. I am a Ph.D.\ candidate in Economics at Johns Hopkins University, advised by Professors Christopher D.~Carroll, Nicholas W.~Papageorge, and Stelios Fourakis, and I expect to complete my degree in May~2026. My research lies at the intersection of macroeconomics, applied microeconomics, and wealth inequality, with an emphasis on how heterogeneity in returns, income risk, and institutions shapes aggregate outcomes and policy transmission.

My job market paper estimates a distribution of heterogeneous returns to wealth using micro data from the Survey of Consumer Finances and embeds these estimates into a structural heterogeneous-agent macroeconomic model. I show that persistent return heterogeneity plays a central role in generating observed wealth concentration and has important implications for macroeconomic policy, particularly for the distributional effects of monetary and fiscal interventions. A related project uses longitudinal data from the Health and Retirement Study to study how trust and beliefs interact with portfolio choices and realized returns, providing a behavioral channel linking micro-level heterogeneity to macroeconomic outcomes. Together, these projects form a coherent research agenda focused on distributional macroeconomics, grounded in micro data and quantitative modeling.

The Georgetown Center for Business and Public Policy is an especially strong fit for my work. The Center’s focus on competition, regulation, and innovation in emerging business models closely aligns with my interest in how institutional environments and regulatory frameworks shape economic behavior and welfare. In particular, the Center’s research on platform markets, innovation incentives, and regulatory design connects naturally to my training in applied microeconometrics and structural modeling. I am interested in extending my research to study how regulation and market structure affect firm entry, pricing behavior, and the distributional consequences of innovation, especially in sectors characterized by rapid technological change.

The Center’s location in Washington, DC, and its active engagement with policymakers and practitioners are also highly attractive to me. I am eager to work in an environment where rigorous economic research directly informs policy discussions and where collaboration across economics, business, and public policy is encouraged. I would welcome the opportunity to present my work in the Center’s seminars, receive feedback from affiliated faculty, and develop new collaborative projects during the fellowship.

Thank you very much for your time and consideration. I would welcome the opportunity to discuss my application further.

\closing{Sincerely,}

\end{letter}

\end{document}
